\notechapter{N-20230117}{More Function Exercises: Periodicity, Piecewise Definitions, Images, and Functional Equations}
\enlargethispage{4\baselineskip}

\section{Piecewise Maps and Their Domains}

Function exercises of this kind are solved by translating each problem into one of several languages:
\[
  \text{formula},\qquad \text{graph},\qquad \text{image/preimage},\qquad \text{iteration},\qquad \text{functional equation}.
\]
Many errors come from using a formula while ignoring its domain or treating an identity valid for all inputs as though it were a formula at a few special values. We follow the standard real-analysis approach to continuous images and the Cauchy equation, as in Gressman's \emph{Advanced Analysis} notes and the lecture notes on Cauchy's functional equation.\footnote{\href{https://www2.math.upenn.edu/~gressman/analysis/04-ivt.html}{P. Gressman's intermediate-value theorem notes}, with complementary calculus exercises from Carleton University.}

A piecewise-defined function may use different formulas on intervals such as
\[
  x<0,\qquad x=0,\qquad x>0.
\]
When studying continuity or solving equations involving such a function, one must first decide which case the input belongs to.  If the input itself is an expression such as $f(x)$ or $x+1$, the relevant case may change during the calculation.

The safe method is:
\begin{enumerate}[(i)]
  \item identify the possible regions of the input;
  \item solve the equation in each region;
  \item check that the solution actually belongs to the assumed region.

\end{enumerate}

The pages include small sketches of graphs.  The standard transformations are:
\[
  y=f(x-a) \quad \text{shift right by }a,
\]
\[
  y=f(x)+b \quad \text{shift up by }b,
\]
\[
  y=-f(x) \quad \text{reflect across the }x\text{-axis},
\]
\[
  y=f(-x) \quad \text{reflect across the }y\text{-axis}.
\]
Absolute values fold graphs:
\[
  y=|f(x)|
\]
reflects the negative part upward, while
\[
  y=f(|x|)
\]
copies the right half to the left.

A piecewise function is a function defined after partitioning the domain into regions.  The mathematical issue is not only which formula applies, but how the pieces meet.  Discontinuities occur at boundaries of the partition; differentiability requires matching first-order behavior across boundaries.

For example, if
\[
  f(x)=\begin{cases}f_1(x),&x<a,\\ f_2(x),&x\ge a,
  \end{cases}
\]
then continuity at $a$ requires
\[
  \lim_{x\to a^-}f_1(x)=f_2(a).
\]
Differentiability additionally requires matching slopes.  This is a local gluing condition. For example, let \(f(x)=ax^2\) for \(x<1\) and \(f(x)=bx+c\) for \(x\ge1\). Continuity at \(1\) requires \(a=b+c\), while differentiability requires in addition \(2a=b\). Hence differentiability holds exactly when \(b=2a\) and \(c=-a\). The boundary test is an actual linear system for the unknown coefficients, rather than an informal statement that the formulas should match.

\section{Images of intervals and finite maps}

For a continuous function on an interval, the image is also an interval.  Indeed, intervals are the connected subsets of \(\mathbb R\), and continuous images of connected sets remain connected. This is the intermediate value theorem in set language. However, the image of an open interval need not be open: \(f(x)=x^2\) sends \((-1,1)\) to \([0,1)\).  If the function is monotone on the interval, then the endpoints determine the image directly.

To find the range of an expression, rewrite it as a one-variable function, determine monotonicity or extrema, and translate the result back into an interval.

The printed exercise page also touches finite functions.  If $|A|=m$ and $|B|=n$, then the number of functions $A\to B$ is
\[
  n^m,
\]
because each element of $A$ independently chooses one image in $B$.

If one requires injectivity and $m\le n$, the number is
\[
  n(n-1)\cdots(n-m+1).
\]
If one requires bijectivity and $m=n$, the number is
\[
  n!.
\]
If surjectivity is required for \(m\ge n\), inclusion--exclusion gives
\[
 \#\operatorname{Surj}(A,B)=\sum_{j=0}^{n}(-1)^j\binom nj(n-j)^m.
\]
Indeed, count all \(n^m\) functions and subtract those omitting at least one of the \(n\) target elements. This explains why injectivity and surjectivity lead to different counting problems.

\section{Functional equations as transformations}

A relation such as
\[
  f(x+1)=1-f(x)
\]
implies
\[
  f(x+2)=1-f(x+1)=f(x).
\]
Thus the function is periodic with period $2$.  The general method is to apply the same equation repeatedly until the input returns to a familiar form.

If instead
\[
  f(x+T)=f(x)
\]
for all $x$, then $T$ is a period.  The smallest positive period, if it exists, is called the fundamental period.

Substituting special values such as $x=0$, $y=0$, or $y=x$ is a standard strategy.  For a functional equation in two variables,
\[
  F(x,y,f(x),f(y),f(x+y))=0,
\]
choosing simple values isolates pieces of the function.

For example, if
\[
  f(x+y)=f(x)+f(y)+c,
\]
then setting $x=y=0$ determines $f(0)$, and subtracting a constant can reduce the equation to the additive Cauchy equation.

A functional equation often specifies a transformation rule rather than a direct formula.  For example, equations involving $f(x+1)$, $f(2x)$, or $f(f(x))$ are naturally studied by iteration.  The orbit of $x$ under a transformation $T$ is
\[
  x,\ T(x),\ T^2(x),\ \ldots.
\]
A functional equation may prescribe how $f$ behaves along these orbits.

For instance, if
\[
  f(x+1)=f(x),
\]
then $f$ is constant along orbits of the translation $T(x)=x+1$.  Periodicity is invariance under a group action of $\mathbb Z$.

A functional equation often says that $f$ is invariant or equivariant under transformations.  The equation
\[
  f(x+T)=f(x)
\]
says $f$ is invariant under the translation action of $T\mathbb Z$.  The equation
\[
  f(ax)=bf(x)
\]
says $f$ is homogeneous with respect to scaling.

Thus solving a functional equation often means finding the invariants of a group or semigroup action.  This perspective makes the original examples more natural: periodicity, parity, and scaling are all symmetry constraints.

\section{Regularity makes Cauchy's equation rigid}

Cauchy's equation is
\[
  f(x+y)=f(x)+f(y).
\]
For an additive map \(f:\mathbb R\to\mathbb R\), continuity at one point, Lebesgue measurability, or boundedness on a nondegenerate interval each implies
\[
  f(x)=cx.
\]
Here is the continuous case. Additivity gives \(f(0)=0\), \(f(nx)=nf(x)\), and \(f(p/q)=(p/q)f(1)\) for rationals \(p/q\). If \(q_k\in\mathbb Q\) approaches real \(x\), continuity yields \(f(x)=\lim_k f(q_k)=x f(1)\). The discontinuous Hamel-basis solutions show exactly what continuity rules out.
But without regularity assumptions, there are wild solutions constructed using a Hamel basis of $\mathbb R$ over $\mathbb Q$.

This shows why elementary functional equations often include hidden conditions such as continuity or monotonicity.  These conditions are not technical decorations; they exclude pathological solutions arising from the axiom of choice.

\section{From iterated intervals to symbolic dynamics}

Piecewise-defined functions and iteration can lead to symbolic dynamics.  If an interval map sends subintervals across one another, one can encode an orbit by a sequence of symbols recording which subinterval contains each iterate.

The shift map
\[
  \sigma((a_0,a_1,a_2,\ldots))=(a_1,a_2,a_3,\ldots)
\]
is a central model.  Complicated interval dynamics can often be compared to shifts on symbol spaces.

For a specific example, partition \([0,1)\) into \(I_0=[0,1/2)\) and \(I_1=[1/2,1)\), and set \(T(x)=2x\pmod1\). The itinerary of \(x\) is the sequence \(a_k\in\{0,1\}\) recording which interval contains \(T^k(x)\). Under the binary expansion convention that avoids trailing strings of ones, this itinerary is precisely the binary digits of \(x\); the map \(T\) shifts them left. Thus symbolic dynamics is the exact encoding of a concrete interval map, with the usual care about dyadic endpoints.
